% --- 题目 ---
\problem[P131-1 (25-1)]{已知级数：\textcircled{1} $\displaystyle \sum_{n=1}^{\infty} \sin \frac{n^3 \pi}{n^2+1}$； \textcircled{2} $\displaystyle \sum_{n=1}^{\infty} (-1)^n \left(\frac{1}{\sqrt[3]{n^2}} - \tan \frac{1}{\sqrt[3]{n^2}}\right)$, 则 ( \quad )
\begin{tasks}(2)
  \task \textcircled{1}与\textcircled{2}均条件收敛.
  \task \textcircled{1}条件收敛, \textcircled{2}绝对收敛.
  \task \textcircled{1}绝对收敛, \textcircled{2}条件收敛.
  \task \textcircled{1}与\textcircled{2}均绝对收敛.
\end{tasks}}
\ansat{315；【十年真题】 - 考点：常数项级数的收敛性 - 1}

% --- 题目 ---
\problem[P131-3 (23-1,3)]{已知 $\displaystyle a_n < b_n$ \ ($\displaystyle n = 1, 2, \dots$). 若级数 $\displaystyle \sum_{n=1}^{\infty} a_n$ 与 $\displaystyle \sum_{n=1}^{\infty} b_n$ 均收敛, 则 “$\displaystyle \sum_{n=1}^{\infty} a_n$ 绝对收敛” 是 “$\displaystyle \sum_{n=1}^{\infty} b_n$ 绝对收敛” 的 ( \quad )}
\begin{tasks}(2)
  \task 充分必要条件.
  \task 充分不必要条件.
  \task 必要不充分条件.
  \task 既不充分也不必要条件.
\end{tasks}
\ansat{315；【十年真题】 - 考点：常数项级数的收敛性 - 3}

% --- 题目 ---
\problem[P131-4 (19-1)]{设 $\displaystyle \{u_n\}$ 是单调增加的有界数列, 则下列级数中收敛的是 ( \quad )
\begin{tasks}(4)
  \task $\displaystyle \sum_{n=1}^{\infty} \frac{u_n}{n}.$
  \task $\displaystyle \sum_{n=1}^{\infty} (-1)^n \frac{1}{u_n}.$
  \task $\displaystyle \sum_{n=1}^{\infty} \left(1 - \frac{u_n}{u_{n+1}}\right).$
  \task $\displaystyle \sum_{n=1}^{\infty} \left( u_{n+1}^2 - u_n^2 \right).$
\end{tasks}}
\ansat{315；【十年真题】 - 考点：常数项级数的收敛性 - 4}

% --- 题目 ---
\problem[P131-5 (19-3)]{若 $\displaystyle \sum_{n=1}^{\infty} n u_n$ 绝对收敛, $\displaystyle \sum_{n=1}^{\infty} \frac{v_n}{n}$ 条件收敛, 则 ( \quad )
\begin{tasks}(2)
  \task $\displaystyle \sum_{n=1}^{\infty} u_n v_n$ 条件收敛.
  \task $\displaystyle \sum_{n=1}^{\infty} u_n v_n$ 绝对收敛.
  \task $\displaystyle \sum_{n=1}^{\infty} \left( u_n + v_n \right)$ 收敛.
  \task $\displaystyle \sum_{n=1}^{\infty} \left( u_n + v_n \right)$ 发散.
\end{tasks}}
\ansat{315；【十年真题】 - 考点：常数项级数的收敛性 - 5}

% --- 题目 ---
\problem[P131-6 (17-3)]{若级数 $\displaystyle \sum_{n=2}^{\infty} \left[\sin \frac{1}{n} - k \ln\left(1 - \frac{1}{n}\right)\right]$ 收敛, 则 $\displaystyle k = $ ( \quad )
\begin{tasks}(4)
  \task $\displaystyle 1$.
  \task $\displaystyle 2$.
  \task $\displaystyle -1$.
  \task $\displaystyle -2$.
\end{tasks}}
\ansat{315；【十年真题】 - 考点：常数项级数的收敛性 - 6}

% --- 题目 ---
\problem[P131-8 (16-1)]{已知函数 $\displaystyle f(x)$ 可导, 且 $\displaystyle f(0)=1, 0 < f'(x) < \frac{1}{2}$.
设数列 $\displaystyle \{x_n\}$ 满足 $\displaystyle x_{n+1} = f(x_n) \ (n=1, 2, \dots)$. 证明:
\begin{enumerate}
\item[(1)] 级数 $\displaystyle \sum_{n=1}^{\infty} (x_{n+1} - x_n)$ 绝对收敛;
\item[(2)] $\displaystyle \lim_{n \to \infty} x_n$ 存在, 且 $\displaystyle 0 < \lim_{n \to \infty} x_n < 2$.
\end{enumerate}}
\ansat{316；【十年真题】 - 考点：常数项级数的收敛性 - 8}

